% MIN 15:00 à 16:00. Partir des reproductions menées de zéro.
\begin{frame}{La simulation prend en charge la boucle d'entraînement}
    En reproduisant des expériences à partir de zéro, j'ai dû réécrire
    la même boucle d'entraînement.

    \medskip
    \textbf{L'utilisateur} choisit la répartition des données entre participants.

    \begin{block}{La simulation automatise le protocole}
        \begin{enumerate}
            \item Former et tirer les petits lots de données.
            \item Calculer les gradients locaux.
            \item Appliquer l'attaque et l'agrégation.
            \item Mettre à jour le modèle.
        \end{enumerate}
    \end{block}
\end{frame}

% MIN 16:00 à 17:00. Définir la fidélité d'une simulation.
\begin{frame}{Fidélité des simulations}
    \begin{block}{L'objectif}
        La simulation permet de faire varier les paramètres de l'expérience tout en gardant
        la même boucle d'entraînement sous PyTorch.
    \end{block}
    \medskip
    \begin{block}{Fidélité au protocole}
        Une simulation fidèle conserve les mécanismes du protocole nécessaires
        à l'étude du phénomène visé et documente les adaptations apportées.
    \end{block}
\end{frame}
